\chapter{Scripts}\label{ch:scripts}
En este capítulo se presentan los scripts utilizados para la descarga y manejo de los datos necesarios para la simulación.
\begin{lstlisting}[language=bash,label={lst:script_descargar_distances},caption=script bash para realizar la descarga de distancias por partes]
#!/bin/bash
BASE_URL="https://tapvizier.cds.unistra.fr/TAPVizieR/tap/sync"
TABLE="I/352/gedr3dis"
FIELDS="Source,rgeo"
N_PARTS=500

# Define los límites del source_id
SOURCE_MIN=4295806720
SOURCE_MAX=6917528997577384320
STEP=$(( (SOURCE_MAX - SOURCE_MIN) / N_PARTS ))

for ((k = 0; k < N_PARTS; k++)); do
  START=$(( SOURCE_MIN + k * STEP ))
  END=$(( START + STEP - 1 ))
  FILENAME="gedr3dis_range_${k}.csv"

  if [[ -s "$FILENAME" ]]; then
    echo "Parte $k ya existe. Se omite."
    continue
  fi

  echo "Descargando parte $k: source_id entre $START y $END"

  QUERY="SELECT ${FIELDS} FROM \"${TABLE}\" WHERE Source BETWEEN ${START} AND ${END}"
  ENCODED_QUERY=$(echo "$QUERY" | sed 's/ /+/g' | sed 's/"/%22/g')
  URL="${BASE_URL}?REQUEST=doQuery& \
      LANG=ADQL&FORMAT=csv&QUERY=${ENCODED_QUERY}"

  for attempt in {1..3}; do
    echo "Intento $attempt..."
    wget -O "$FILENAME" "$URL" && break
    echo "Falló el intento $attempt para parte $k"
    sleep 5
  done

  if [[ ! -s "$FILENAME" ]]; then
    echo "Falló la descarga de parte $k tras 3 intentos." >&2
    rm -f "$FILENAME"
  else
    echo "Parte $k descargada correctamente."
  fi
done
\end{lstlisting}
\begin{lstlisting}[language=Awk,label={lst:script_reducir_astro},caption=script awk para reducir astro de forma paralela]
#!/bin/bash

# Crear carpeta de salida si no existe
mkdir -p temp
mkdir -p ../reducidos_astro

# Función que procesa un solo archivo
procesar_archivo() {
    local file="$1"
    local index="$2"
    local clean_index=$((10#$index))
    local temp_csv="temp/temp_${clean_index}.csv"
    local reduced_file="../reducidos_astro/astro$(printf "%02d" "$clean_index").csv"

    echo "Procesando archivo: $file..."

    # Descomprimir sin borrar el .gz
    gunzip -c "$file" > "$temp_csv"

    # Procesar con awk y guardar el archivo reducido
    awk -F',' '
    BEGIN { foundHeader = 0 }
    /^#/ { next }
    !foundHeader {
        foundHeader = 1
        for (i=1; i<=NF; i++) colName[$i] = i
        print $colName["source_id"], $colName["mass_flame"], $colName["radius_flame"], $colName["logg_gspphot"]
        next
    }
    {
        print $colName["source_id"], $colName["mass_flame"], $colName["radius_flame"], $colName["logg_gspphot"]
    }' OFS=',' "$temp_csv" > "$reduced_file"

    rm -f "$temp_csv"
}

export -f procesar_archivo

# Obtener lista de archivos y numerarlos
ls *.gz | nl -v 1 -n rz | while read -r index file; do
    echo "$file $index"
done > archivos_con_indices.txt

# Ejecutar en paralelo
cat archivos_con_indices.txt | parallel -j 16 --colsep ' ' procesar_archivo {1} {2}

echo "Proceso completado. Archivos reducidos en la carpeta 'reducidos_astro/'."
\end{lstlisting}
\begin{lstlisting}[language=Python,label={lst:script_combinar},caption=script para combinar todas las tablas]
import pandas as pd
import duckdb
import glob
import os
from concurrent.futures import ProcessPoolExecutor
import time

# Configuración de paths
astro_folder = 'reducidos_astro'
mainsource_folder = 'reducidos'
distance_folder = 'distances'
output_folder = 'optimizados'
db_path = 'distance_data.duckdb'

# Crear carpeta de salida
os.makedirs(output_folder, exist_ok=True)

def crear_base_duckdb(distance_folder, db_path):
    print("Creando base DuckDB desde archivos de distancia...")
    distance_files = sorted(glob.glob(os.path.join(distance_folder, '*.csv')))
    if not distance_files:
        raise FileNotFoundError("No se encontraron archivos en 'distances'.")

    chunks = []
    for i, file in enumerate(distance_files):
        print(f"Leyendo archivo {i+1}/{len(distance_files)}: {file}")
        df = pd.read_csv(file)
        chunks.append(df)

    distance_df = pd.concat(chunks, ignore_index=True)
    if 'Source' in distance_df.columns:
        print("Renombrando columna 'Source' a 'source_id'")
        distance_df.rename(columns={'Source': 'source_id'}, inplace=True)

    con = duckdb.connect(db_path)
    con.execute("DROP TABLE IF EXISTS distance")
    con.register("df", distance_df)
    con.execute("CREATE TABLE distance AS SELECT * FROM df")
    con.close()
    print("Base de datos DuckDB creada")

def merge_and_save(i, main_file, astro_file, db_path, output_folder):
    try:
        print(f"[{i:04d}] Procesando archivos: {os.path.basename(main_file)}, {os.path.basename(astro_file)}")
        main_df = pd.read_csv(main_file)
        astro_df = pd.read_csv(astro_file)

        if 'parallax' in main_df.columns:
            main_df.drop(columns=['parallax'], inplace=True)

        merged = pd.merge(main_df, astro_df, on='source_id', how='inner')

        source_ids = tuple(merged['source_id'].unique())
        if len(source_ids) == 1:
            source_ids = f"({source_ids[0]})"
        else:
            source_ids = str(source_ids)

        con = duckdb.connect(db_path, read_only=True)
        distance_df = con.execute(
            f"SELECT * FROM distance WHERE source_id IN {source_ids}"
        ).fetchdf()
        con.close()

        merged = pd.merge(merged, distance_df, on='source_id', how='left')

        output_file = os.path.join(output_folder, f'merged_{i:04d}.csv')
        merged.to_csv(output_file, index=False)
        return f'Guardado: {output_file}'

    except Exception as e:
        return f' {i:04d}: {str(e)}'

def run():
    start_time = time.time()
    print("Iniciando procesamiento completo...")

    # Paso 1: Crear DB
    crear_base_duckdb(distance_folder, db_path)

    # Paso 2: Preparar listas de archivos
    main_files = sorted(glob.glob(os.path.join(mainsource_folder, '*.csv')))
    astro_files = sorted(glob.glob(os.path.join(astro_folder, '*.csv')))

    if len(main_files) != len(astro_files):
        raise ValueError("Diferente número de archivos en main y astro")

    print(f"Total de pares a procesar: {len(main_files)}")

    # Paso 3: Merge en paralelo
    with ProcessPoolExecutor(max_workers=16) as executor:
        tasks = [
            executor.submit(merge_and_save, i, main_file, astro_file, db_path, output_folder)
            for i, (main_file, astro_file) in enumerate(zip(main_files, astro_files))
        ]
        for future in tasks:
            print(future.result())

    # Paso 4: Eliminar base de datos temporal
    if os.path.exists(db_path):
        os.remove(db_path)
        print(f" Base de datos eliminada: {db_path}")

    print(f" Proceso finalizado en {time.time() - start_time:.2f} segundos")

if __name__ == '__main__':
    run()
\end{lstlisting}
